\section{The Case II rank-four support calculation}
\label{sec:sr-rank-four-support}
The specialization in Section~\ref{sec:sr-character-integers} eliminates
the other rank-four cases. We give the remaining core calculation
from its stated tower, common-space and modular-support inputs.
The existence of the linked automorphic seed and the finite-level
automorphic modules remains explicit.

\subsection{The rank-four characteristic-\(\ell\) support calculation}
\label{subsec:sr-rank-four-core}
This subsection supplies the rank-four inputs to the common-space
completion and component-propagation lemmas. It is restricted to the
actual Case II tower and requires no Case III construction.

Write \(F\) for its totally real abelian field,
\(d=[F:\mathbb Q]\), \(s=\#\{v\mid\ell\}\), and
\(\delta=[F_v:\mathbb Q_\ell]\), so \(s\delta=d\).
Let \(\Lambda\) be a finite coefficient DVR with residue field \(k\)
of odd characteristic \(\ell\).
The prescribed residual characters are \(\chi_0,\ldots,\chi_3\),
with \(\chi_i\chi_{3-i}=\nu\), pairwise distinct at \(g_0\).
Representations have fixed multiplier \(\nu\), are diagonal at
\(g_0\), satisfy the prescribed global ramification condition,
and have ordinary flags represented in the lower open cell by
lower unipotent symplectic matrices \(n_v\).
Their ordinary characters are trivial on the maximal
\(\ell\)-power torsion subgroup of the local abelianization.
The tower makes the residual representation trivial on inertia
at \(\ell\); its multiplier is trivial there in characteristic
\(\ell\).

Let \(B\subset A\) be the common space: \(B\) is complete,
\(A\) is of finite type over \(B\), the universal matrix and flag
entries generate \(A\) over \(B\), and \(B\) contains all
characteristic-polynomial coefficients and ordinary diagonal
characters \(\kappa_{v,i}\).
The completion \(R\) at a residual point represents the normalized
deformation problem. The common-space construction and its completion/dimension transport
are retained as hypotheses: the field-transcendence bounds below must
give the indicated complete-local dimension bounds. Their verification
for the actual common space is part of its construction, and is not
inferred just from the existence of a universal formal ring.

The property of a modular ideal \(J\subset B\) needed below is
\begin{equation}\label{eq:sr-pm-finite}
 \mathfrak m_B^e\subset J+\mathfrak m_\Lambda B+
 (\kappa_{v,i}(t)-1:v\mid\ell,\ 0\le i\le3,\ t\in I_v)
 \quad\text{for some }e.
\end{equation}
We also require its actual curve propagation property:
if \(A\to k'[[T]]\), with \(k'\) finite, is centered over
\(\mathfrak m_B\), kills \(J\), has absolutely irreducible
generic representation, and has pairwise distinct ordinary
characters on
\(U_v=I_v\cap\ker\epsilon_\ell\cap\ker\bar r\),
then every prime of \(B\) through which this curve factors
contains \(J\). These are properties of an actual ideal; no
existence of such an ideal is asserted in this subsection.

\begin{srlemma}[Rank-four trace fiber]\label{lem:sr-four-trace-fiber}
At an absolutely irreducible field-valued point, the matrix and
flag entries are algebraic over the image field of \(B\) and
at most \(2+4s\) additional elements. Consequently, for every
irreducible generic prime \(\mathfrak q\) of the representing
completion,
\[
 \dim R/\mathfrak q\le
       \dim B/\mathfrak q_B+2+4s.
\]
\end{srlemma}
\begin{proof}
Pass to an algebraic closure of the point field and put
\(r(g_0)=\operatorname{diag}(d_0,d_1,d_2,d_3)\).
These distinct eigenvalues are algebraic over the image of \(B\).
Let \(E\) be an algebraically closed intermediate field containing
them, all traces, and all multiplier values.

The four identities
\[
 \operatorname{tr}(r(g_0)^m r(h))
       =\sum_i d_i^m r(h)_{ii}\qquad(0\le m\le3)
\]
and Vandermonde inversion put every diagonal entry in \(E\).
Apply the same inversion to the \(00\) entry of
\(r(u)r(g_0)^m r(w)\), and then twice to the \(00\) entry of
\(r(u)r(g_0)^m r(g)r(g_0)^n r(w)\). They give
\begin{equation}\label{eq:sr-four-triples}
 r(u)_{0i}r(w)_{i0}\in E,\qquad
 r(u)_{0i}r(g)_{ij}r(w)_{j0}\in E .
\end{equation}
The matrix-algebra spanning theorem supplies
\(a_i=r(u_i)_{0i}\ne0\) and \(b_i=r(w_i)_{i0}\ne0\).
Adjoin only \(a_1,a_2\); \(a_0\) is already in \(E\).

Our adjoint convention is
\(X^\#_{ij}=\sigma_i\sigma_jX_{3-j,3-i}\),
\(\sigma=(1,1,-1,-1)\).
The symplectic relation gives
\(-r(g)_{02}=\nu(g)r(g^{-1})_{13}\).
Thus \(c=r(u_2^{-1})_{13}=-a_2/\nu(u_2)\) is nonzero in \(E\).
Now \(a_1cb_3\in E\) gives \(b_3\in E\);
\(a_3b_3\in E\) gives \(a_3\in E\).
Then \(a_ib_i\in E\) gives every \(b_i\), and the triple
products give every entry of \(r(g)\).

Every lower unipotent symplectic flag matrix is
\[
 \begin{pmatrix}
 1&0&0&0\\ A&1&0&0\\ B&D&1&0\\
 C&B-DA&-A&1
 \end{pmatrix}.
\]
Four additional coordinates per place therefore suffice.
Coordinate generation and the common-space dimension-transport
lemma prove the result.
\end{proof}

\begin{srlemma}[Reducible tangent bound]
\label{lem:sr-four-reducible-tangent}
Suppose the tower cohomology bounds give
\[
 s\le C,\quad \dim Z_w\le C\ (w=0,\psi,-\psi),\quad
 \dim Z_w\le d+C\ (w\text{ odd}),\quad
 \dim S_w\le C\ (w\text{ a root}),
\]
where \(Z_w\) denotes the globally allowed normalized cocycles
and \(S_w\subset Z_w\) the subspace locally a coboundary.
At a good residual point every proper reducible support quotient
has tangent dimension at most \(d+20C\) and dimension at most
\(d+20C+1\).
\end{srlemma}
\begin{proof}
Put
\[
 \lambda=(1,0,1,0),\quad\mu=(1,1,0,0),\quad
 w_{ij}=(\lambda_i-\lambda_j,\mu_i-\mu_j),
\]
and \(\alpha=(1,0),\beta=(-1,1),\psi=(0,1),\theta=(1,1)\).
The odd roots are \(\pm\alpha,\pm\beta,\pm\theta\).
A support \(S\) means \(r(g)_{ij}=0\) when
\(i\ne j,w_{ij}\notin S\).
The nine maximal proper supports and the pivots needed below are
\[
\begin{array}{c|l|c|c}
&S&\text{split}&\text{partially split}\\\hline
1&\{\beta,-\beta,-\alpha,-\theta,-\psi\}&12&12\\
2&\{\beta,-\beta,\alpha,\theta,\psi\}&-&-\\
3&\{\theta,-\theta,\psi,\beta,-\alpha\}&-&12\\
4&\{\theta,-\theta,-\psi,-\beta,\alpha\}&-&-\\
5&\{\alpha,-\alpha,-\theta,-\psi,-\beta\}&01&01\\
6&\{\alpha,-\alpha,\theta,\psi,\beta\}&-&01\\
7&\{\psi,-\psi,\beta,-\alpha,-\theta\}&12&12\\
8&\{\psi,-\psi,-\beta,\alpha,\theta\}&-&-\\
9&\{\beta,-\beta,\theta,-\theta\}&-&03
\end{array}
\]
This list covers reducibility because an invariant subspace for
the regular diagonal \(r(g_0)\) is a coordinate subspace.
There are four line stabilizers, four coordinate Lagrangian
plane stabilizers, and one pair of complementary nondegenerate
coordinate planes; three-dimensional subspaces are handled by
their symplectic perpendicular lines.

The split alternative applies when
\[
 (-\psi\in S\text{ and }(\beta\in S\text{ or }-\alpha\in S))
 \quad\text{or}\quad\{-\alpha,-\beta\}\subset S.
\]
The partially split alternative applies when
\(-\alpha\in S\) or \(\{\beta,-\theta\}\subset S\).
Goodness of the residual point guarantees one of these
alternatives, and in the second it gives the partially split
shape described next. Thus rows \(2,4,8\) cannot occur in the
corresponding alternative.

At a split point \(r=D=\operatorname{diag}(\chi_i)\), \(n_v=1\).
For a tangent vector \((z,(Y_v))\), each \(z_{ij}\) belongs to
\(Z_{w_{ij}}\). The linearized ordinary condition makes every
lower entry locally a coboundary, and off-support entries vanish.
Among the independent upper odd entries \(01,03,12\), each
applicable row allows only the displayed pivot. Killing it
therefore makes all independent odd entries strict.

At a partially split point \(r=D\) and
\(n_v=1+b_vX_-\), where \(b_v\ne0\) and
\(X_-=E_{20}+E_{31}\).
The residual flag equation, in its \(20\) entry, gives
\(\chi_0|_{G_v}=\chi_2|_{G_v}\).
Thus roots with the same first coordinate induce the same local
character. Direct multiplication in \(\mathfrak{sp}_4\) gives
\[
 Z_{10}=-b_vZ_{12},\quad Z_{30}=-b_v^2Z_{12},\quad
 Z_{21}=2b_vZ_{01}+b_v^2Z_{03}
\]
whenever \(n_v^\#Zn_v\) is upper triangular.
Apply this to
\(Z=z(s)+D(s)Y_vD(s)^{-1}-Y_v\), and pass modulo local
coboundaries. We obtain
\[
 [z_{12}]=0\Longleftrightarrow[z_{10}]=0
            \Longleftrightarrow[z_{30}]=0,\qquad
 [z_{21}]=2b_v[z_{01}]+b_v^2[z_{03}].
\]
Both \(2b_v\) and \(b_v^2\) are invertible.
In each indicated row, its pivot together with off-support
zeros therefore kills all six independent odd classes.
This proves the finite table rather than assuming it.

The pivot image has dimension at most \(d+C\).
On its kernel the ten independent symplectic entries
\(i\le3-j\) are strict odd cocycles or even/zero cocycles,
of total dimension at most \(10C\).
The six lower entries of \(n_v^\#Y_vn_v\) per place have
dimension at most \(6s\le6C\).
These coordinates jointly determine the vector, since symplectic
symmetry recovers the remaining global entries and the transformed
local matrices are strictly lower triangular.
Thus its dimension is at most \(d+17C\le d+20C\).

The reduction to the two simple points uses an elementary finite
weight filtration. The residual shape is either supported in
\(\{\beta,-\alpha,-\theta,-\psi\}\), with local flag support
\(\{-\alpha,-\psi,-\theta\}\), or in
\(\{-\alpha,-\beta,-\psi,-\theta\}\).
Use diagonal conjugation weights \((0,2,1,3)\) in the first
case, and \((0,1,2,3)\) in the second.
The limits are \(D,1\).
For the partially split shape the global support is
\(\{\beta,-\alpha,-\theta\}\), the same local flag support
has \(n_{v,20}\ne0\), and \((0,1,0,1)\) gives
\(D,1+n_{v,20}X_-\).
Filter tangent matrices by their entry weights \(m_i-m_j\).
Taking the lowest nonzero homogeneous part preserves every
linearized cocycle, support, symplectic and flag equation:
all structure matrices have nonnegative degree and the limits
are their degree-zero parts.
The associated graded embeds, degree by degree, into the tangent
space of the stated limit. There are finitely many weights,
and the degree images are independent, proving the dimension
inequality without a smoothness assertion.

Finally a derivation of the representing support quotient
injects into this matrix tangent space. If its derivatives on
all universal representation and flag entries vanish, its two
maps to \(k[\varepsilon]/\varepsilon^2\) give the same universal
deformation and hence coincide.
Nakayama then generates the maximal ideal by a uniformizer and
at most \(d+20C\) further elements, proving the dimension bound.
\end{proof}

\begin{srlemma}[Ordinary-character collision]\label{lem:sr-four-collision}
Assume \eqref{eq:sr-pm-finite}.
Let \(\mathfrak q\subset B\) be a prime containing
\(J+\mathfrak m_\Lambda B\). If
\(\kappa_{v,i}=\kappa_{v,j}\) on \(U_v\) modulo \(\mathfrak q\),
for some \(i\ne j\), then
\[
 \dim B/\mathfrak q+\delta\le2d+1.
\]
\end{srlemma}
\begin{proof}
In a local domain of characteristic \(\ell\),
\(1+\mathfrak m\) has no nontrivial finite-order elements.
For \(\ell\)-power order use
\((x-1)^{\ell^r}=x^{\ell^r}-1\).
For prime-to-\(\ell\) order \(n\), the geometric sum is a
unit because its residue is \(n\); then separate these factors
for general finite order.

Local class field theory now gives \(\delta\) pinning elements
of \(I_u\) at each \(u\mid\ell\): a continuous
\(G_u\to1+\mathfrak m\) character trivial at those elements is
trivial on \(I_u\).
Indeed the torsion-free pro-\(\ell\) quotient of the local unit
group is \(\mathbb Z_\ell^\delta\), and its finite torsion is
killed by the preceding observation.
Also choose \(t_0\in I_v\) such that
\(\{t_0\}\cup(I_v\cap\ker\epsilon_\ell)\) pins such characters
on \(I_v\).
The cyclotomic image is a closed subgroup of
\(1+\ell\mathbb Z_\ell\), hence procyclic.
The extra kernel of \(\bar r\) in \(U_v\) makes no difference,
since the tower kills the residual inertia action.

Choose \(a\) with \(\{a,3-a,i,3-i\}=\{0,1,2,3\}\) and either
\(a=j\) or \(j=3-i\). In \(B/\mathfrak q\) impose
\[
 \begin{array}{ll}
 \kappa_{u,a}(t)-1&
    \text{at all \(\delta\) pins for all \(u\)},\\
 \kappa_{u,i}(t)-1&
    \text{at those pins for all \(u\ne v\)},\\
 \kappa_{v,i}(t_0)-\kappa_{v,j}(t_0).&
 \end{array}
\]
There are \(d+(d-\delta)+1=2d-\delta+1\) parameters,
all in the maximal ideal.
In every prime quotient containing them,
\(\kappa_{u,a}|_{I_u}=1\) for all \(u\), and
\(\kappa_{u,i}|_{I_u}=1\) for \(u\ne v\).
The original collision and the last parameter give equality on
all of \(I_v\).
If \(a=j\), this kills \(\kappa_{v,i}\);
if \(j=3-i\), the pairing
\(\kappa_{v,i}\kappa_{v,3-i}=1\) gives
\(\kappa_{v,i}^2=1\), which again kills it.
Pairing kills the remaining characters on all inertia.
Equation \eqref{eq:sr-pm-finite} forces this prime to be maximal.
Thus the parameter ideal has maximal radical.
The parameter dimension theorem gives
\(\dim B/\mathfrak q\le2d-\delta+1\).
\end{proof}

\begin{srlemma}[Normalized presentation]\label{lem:sr-four-presentation}
Let \(h^0,h^2\) be the degree zero and degree two dimensions of
the ordinary global-to-local adjoint cone, and \(h^1_c\) its
degree one dimension. The arithmetic cone estimate
\[
 h^2+h^0+6d\le h^1_c+4d
\]
gives a presentation
\[
 R=\Lambda[[x_1,\ldots,x_m]]/(f_1,\ldots,f_n),
 \qquad n+2+2d\le m+2s.
\]
\end{srlemma}
\begin{proof}
The fixed-similitude adjoint \(\mathfrak{sp}_4(k)\) has dimension
\(10\); an ordinary symplectic Borel has dimension \(6\);
the diagonal symplectic torus has dimension \(2\).

Temporarily omit the local \(\ell\)-power torsion condition.
For a small extension with square-zero kernel, choose lifts of
the representation and flags.
Their multiplication and flag defects form respectively a
global \(2\)-cochain and local \(1\)-cochains modulo the Borel.
Associativity is exactly the cone \(2\)-cocycle identity.
Changing lifts changes the defects by its coboundary.
Conversely a coboundary corrects both defects by its global
\(1\)-cochain and local \(0\)-cochains.
All corrections lie in \(\mathfrak{sp}_4\), maintaining the
multiplier; \(2\) is invertible.
Thus cone \(H^2\) is an obstruction space.
The minimal-relator argument gives at most \(h^2\) relations:
minimal relators map injectively to their universal lifting
obstructions, since a zero obstruction corrects the universal
lift and removes the proposed relation.

The normalized tangent dimension is at least \(h^1_c+2-h^0\).
To check the count, represent ordinary degree one classes by
global cocycles and local corrections chosen in the lower
complement of each Borel.
The global conjugation parameters have dimension \(10\), with
kernel of dimension at most \(h^0\).
Normalizing at \(g_0\) imposes at most \(10-2=8\) conditions.
The resulting dimension inequality is
\(h^1_c+2\le\dim T+h^0\).
Dual numbers identify \(T\) with the tangent space of the
representing ring.

Restore the local torsion condition. A generator of the
\(\ell\)-power torsion in each local abelianization suffices.
Its multiplier is \(1\): the cyclotomic image has no
\(\ell\)-power torsion for odd \(\ell\), and the remaining
finite multiplier has order prime to \(\ell\), including the
finite \(2\)-power character in this packet.
Only the first two of the four paired ordinary characters need
to be set equal to \(1\). This adds \(2s\) equations.
Therefore
\[
 n\le h^2+2s,\qquad m\ge h^1_c+2-h^0,
\]
and the cone estimate gives \(m-n\ge2d+2-2s\).

For this cone estimate the local Borel dimension is \(6\),
and the complex-conjugation fixed adjoint dimension is \(4\):
an anti-symplectic involution has two isotropic eigenspaces of
dimension \(2\), and its fixed symplectic Lie algebra is
\(\mathfrak{gl}_2\).
The required absence of a cyclotomic quotient is checked on the
residual adjoint root characters. The actual tower supplies,
for each such character, an element where it differs from the
mod-\(\ell\) cyclotomic character; the zero character is also
separated by complex conjugation. This is exactly the
noncyclotomic condition in the tower construction, rather than
an inference from pairwise distinct values alone.
Thus the arithmetic cone estimate applies to this adjoint.
\end{proof}

\begin{srproposition}[Rank-four component propagation]
\label{prop:sr-four-core}
Assume the common-space construction, the tower cohomology
bounds, \eqref{eq:sr-pm-finite}, and the actual modular curve
propagation property stated above.
If \(d\ge30(C+1)^2\), then at a good residual point, one
component containing \(JA\) implies that every component
through that point contains \(JA\).
\end{srproposition}
\begin{proof}
Put \(N_0=2d+1-2s\).
As \(s\le C\), \(s\delta=d\), and \(d\ge30(C+1)^2\),
\[
 \delta\ge30(C+1),\quad
 d+20C+1<N_0,\quad
 2d+1+2+4s<N_0+\delta.
\]
For the last inequality it suffices that \(2+6s<\delta\).

Lemma~\ref{lem:sr-four-presentation} and the complete-local
component-chain lemma join minimal components by a chain whose
neighbors meet in dimension at least \(N_0+1\).
Lemma~\ref{lem:sr-four-reducible-tangent} makes every reducible
bad set have dimension less than \(N_0\).
Outside these sets, a bad prime where \(J\), the uniformizer,
and an ordinary collision vanish satisfies
\[
 \dim R/\mathfrak q\le
 \dim B/\mathfrak q_B+2+4s
 \le2d+1-\delta+2+4s<N_0
\]
by the trace-fiber and collision lemmas.

Suppose \(J\) vanishes on the first of two neighbors.
The rank-independent prime-avoidance and curve lemma, applied
to their intersection and the finite collection of these bad
ideals, gives a characteristic-\(\ell\) curve with finite
residue field. It kills \(J\), has irreducible generic
representation, and has pairwise distinct ordinary characters
on each \(U_v\).
The modular curve property implies that \(J\) is contained in
the contraction of the second component, since that prime
also lies under the curve kernel.
Thus \(J\) vanishes on the second component.
Induction along the chain proves the assertion in the
completion; faithfully flat transport gives it in \(A\).
\end{proof}

\begin{srlemma}[Characteristic-\(\ell\) consequence]
\label{lem:sr-four-characteristic-ell}
Suppose an actual principal Case II point factors through such
a modular ideal and has characteristic \(\ell\).
Its characteristic polynomial at every \(g\in G_{\mathbb Q}\)
is then the product of the four prescribed residual characters.
It is neither absolutely irreducible of rank \(4\) nor a block
extension with an irreducible rank-\(2\) constituent.
\end{srlemma}
\begin{proof}
Use the stable valuation lattice and the factorization on \(G_F\).
The prescribed ordinary characters become \(1\) on inertia:
the cyclotomic factor reduces to its mod-\(\ell\) value killed
by the tower, and the other finite characters are unramified.
Equation \eqref{eq:sr-pm-finite} therefore kills every element
of \(\mathfrak m_B\), because the target is a domain.
On \(G_F\), all characteristic-polynomial coefficients are the
given finite-field residual constants.

For \(g\in G_{\mathbb Q}\), some positive power \(g^M\) lies in
\(G_F\).
The eigenvalues of \(r(g)^M\), and therefore of \(r(g)\), are
algebraic over \(\mathbb F_\ell\).
Extend the valuation to an algebraic closure of the fraction
field. It is trivial on the algebraic closure of
\(\mathbb F_\ell\), on which reduction is injective.
Thus the characteristic-polynomial coefficients of \(r(g)\)
are constants, and their reduction identifies them with the
given residual polynomial.
Characteristic-polynomial Brauer--Nesbitt identifies its
semisimplification with the four diagonal characters and gives
both exclusions. Complete characteristic polynomials are used,
so this argument also applies when \(\ell=3\).
\end{proof}

\paragraph{Scope and actual remaining producers.}
For the principal packet with \(a\ge2\), the required coefficient
prime is prime to \(N\): in the resonant branch
\(\ell\le2a+1<q,d\), and in the \(3,7\) branch \(q,d>7\).
Thus \(\chi_P=\psi\), and the nontrivial datum is Case II.
The actual host uses the real dyadic cyclotomic subfield;
its dyadic characters are even and its finite character orders
are powers of \(2\).
That makes the prime-to-\(\ell\) finite-character condition in
the tower automatic. The sufficiently large cyclotomic layer
is still necessary for \(d\ge30(C+1)^2\).

The preceding lemmas reconstruct the independent rank-four
characteristic-\(\ell\) core.
The actual Case II endpoint additionally requires the
construction of its linked automorphic rank-four seed and the
construction of its modular ideal with finite-weight, curve,
seed, purity, and automorphy properties.
After those constructions, one chooses a tame Case II tower,
passes to the sufficiently high layer, propagates from the
seed along the residual links, and factors the actual point
through the modular ideal.
The characteristic-\(\ell\) lemma supplies the contradiction;
the actual purity and automorphy properties apply to
characteristic-zero points. The usual deformation endpoint
then gives finite weight zero, no full upper block, and
automorphy at every irreducible point.
No Case III seed or Case III promotion is needed.

% Source alignment, frozen at alpoge-source-7cbf7a0c:
% WeakShimuraTheory/FibreTraces/FibTTrace.lean (fibTorus);
% WeakShimuraTheory/CompactIntegralFields/FibLow.lean (fib);
% WeakShimuraTheory/RedIdeal.lean (redTypes, Bounds, Good, Res);
% WeakShimuraTheory/RedTanFlow.lean (tan_lowest, tan_flow);
% WeakShimuraTheory/CompactIntegralFields/RedTan.lean (redTan);
% WeakShimuraTheory/CompactIntegralFields/HomOfDer.lean (red);
% WeakShimuraTheory/RedDim.lean (redDim);
% WeakShimuraTheory/CompactIntegralFields/ImportsAlt.lean (nn);
% WeakShimuraTheory/AdjointCoefficients/AbsPres.lean;
% WeakShimuraTheory/AdjointCoefficients/MuQuot.lean;
% WeakShimuraTheory/CompactIntegralFields/IncLoc.lean (x1);
% WeakShimuraTheory/CompactIntegralFields/ChCore.lean (ch);
% ModularityArgument/ThresholdArguments/DenseThresholdCases/
% Gsp4PromodularityCaseTwo.lean
%   (proMod_of_pmEx_of_ctx2_of_kPair_of_star);
% PotentialModularityTower.lean;
% ModularityRankFourOfPromodularity.lean.

\subsection{Generic rank-four curves and regular Frobenius cosets}
\label{subsec:sr-four-generic-curves}
We supply the rank-four generic-curve calculation used by
ordinary support patching, including its regular-coset assertion.

\begin{srlemma}[One local element separates all powers]
\label{lem:sr-four-free-local-element}
Let \(k\) be a finite field of characteristic \(\ell\), and
\(\kappa_0,\ldots,\kappa_3:H\to k[[T]]^\times\) be characters.
Let \(U\subset H\) be a subgroup on which all their residual
characters agree.
If every ratio \(\kappa_i/\kappa_j\), \(i\ne j\), is nontrivial
on \(U\), then some \(u\in U\) has all six ratios
\(\kappa_i(u)/\kappa_j(u)\) of infinite order.
\end{srlemma}
\begin{proof}
The group \(1+Tk[[T]]\) is torsion-free:
separate the \(\ell\)-power factor, using
\((x-1)^{\ell^r}=x^{\ell^r}-1\), from the prime-to-\(\ell\)
factor, whose geometric sum is a unit.
Choose finitely many elements of \(U\) witnessing nontriviality
of the six ratios.
Tensor the torsion-free subgroup generated by their ratio
values with \(\mathbb Q\).
For each ratio, products of these elements to integer powers
are trivial only on a proper rational linear subspace of the
exponent lattice. Finitely many proper rational linear
subspaces do not cover that lattice.
A product with exponents outside their union gives \(u\).
Every nontrivial ratio at \(u\) belongs to the torsion-free
group just described and has infinite order.
\end{proof}

\begin{srlemma}[A large residual ratio excludes virtual commutativity]
\label{lem:sr-four-noncommutative}
Let \(\Gamma\) be profinite, \(Z\triangleleft\Gamma\) closed,
and let \(r:\Gamma\to\mathrm{GL}_4(K)\) be continuous and
absolutely irreducible, where \(K\) is an algebraically closed
valued field. Suppose its characteristic polynomials are
integral and reduce to
\(\prod_i(X-\chi_i(g))\) for four residual characters.
Assume some \(\chi_i/\chi_j\) has order greater than \(4\).
If \(Z\) contains \(u\) for which every positive power \(r(u)^m\)
has four distinct eigenvalues, then every open normal subgroup
of \(Z\) has noncommutative image.
\end{srlemma}
\begin{proof}
Suppose \(N\triangleleft Z\) has finite index \(m\) and
commutative image. Its normal core \(C\) in \(\Gamma\)
contains \(u^m\): every conjugate of \(u\) is in \(Z\), whose
\(m\)-th power is in \(N\).
The image of \(C\) commutes and centralizes the regular
semisimple \(r(u^m)\), hence is diagonal in its eigenbasis.
Because \(C\) is normal, \(\Gamma\) permutes its four common
eigenlines. The characters on these lines are distinct, as
they are distinct at \(u^m\).
Absolute irreducibility makes that permutation transitive.
Thus \(r\) is monomial, with permutation map
\(\pi:\Gamma\to S_4\).

For \(g\in\ker\pi\), write its four eigenvalues as \(a_i(g)\).
They are units since the characteristic polynomials of both
\(g\) and \(g^{-1}\) are integral.
If \(h\) sends the \(j\)-th eigenline to the \(i\)-th, the
commutator \(ghg^{-1}h^{-1}\) has eigenvalue
\(a_i(g)/a_j(g)\).
Its residual characteristic polynomial is \((X-1)^4\),
because all \(\chi_i\) are characters.
Hence the four \(a_i(g)\) have the same reduction.
Comparing characteristic polynomials at \(g\) shows
\(\chi_i(g)=\chi_j(g)\) for every \(i,j\).

Thus every residual ratio factors through the subgroup
\(\pi(\Gamma)\subset S_4\).
Its character image is finite cyclic.
Choose an element mapping to a generator of that cyclic image.
An element of \(S_4\) has order \(1,2,3\), or \(4\), so the
ratio has order at most \(4\), a contradiction.
\end{proof}

\begin{srlemma}[Rank-four regular-coset lemma]
\label{lem:sr-four-regular-coset}
Let \(G\) be a profinite group and
\(r:G\to\mathrm{GL}_4(K)\) an absolutely irreducible
representation over an algebraically closed field of
characteristic different from \(2\).
Assume polynomial zero sets of its matrix entries are closed.
Suppose \(r(G)\) contains a semisimple element with all
eigenvalue ratios of infinite order, and every open normal
subgroup of \(G\) has noncommutative image.
Then for every open subgroup \(U\subset G\) and every
\(\sigma\in G\), the coset \(\sigma U\) contains an element
whose characteristic polynomial is separable.
\end{srlemma}
\begin{proof}
Let \(\mathcal G\) be the Zariski closure of \(r(G)\) and
\(\mathcal H=\mathcal G^\circ\).
The preimage of \(\mathcal H\) is closed of finite index,
hence open.
A positive power of the specified semisimple element lies
in \(\mathcal H\), still with all ratios of infinite order.

The restriction of the natural representation to the normal
group \(\mathcal H\) is semisimple by the socle argument:
conjugates of a simple submodule are simple and their sum is
\(\mathcal G\)-invariant.
The isotypic components form one transitive \(\mathcal G\)-orbit.
There is no multiplicity greater than one, since the regular
element in \(\mathcal H\) has no repeated eigenvalues.
Their equal dimensions therefore give one component of
dimension \(4\), two components of dimension \(2\), or four
lines.
The last case would make the open preimage of
\(\mathcal H\) commutative, contrary to hypothesis.

The connected torus which is the identity component of the
closure of the cyclic semisimple subgroup has four distinct
characters in this eigenbasis. Otherwise some eigenvalue
ratio would have finite order.
Choose a cocharacter of that torus with four distinct integer
weights \(w_0,\ldots,w_3\), and write
\(D(z)=\operatorname{diag}(z^{w_i})\).

If \(\mathcal H\) acts irreducibly in dimension \(4\), or the
coset under consideration preserves its two irreducible
\(2\)-blocks, every diagonal entry and every diagonal cofactor
is a nonzero regular function on that coset of \(\mathcal H\).
For a diagonal entry this follows because the orbit of the
corresponding vector spans its irreducible block; for a
cofactor apply the same observation to the inverse matrix.
Connectedness makes \(\mathcal H\) irreducible as a variety.
Thus there is a matrix \(x\) in the coset with all these
entries and cofactors nonzero.

We check directly that \(xD(z)\) has separable characteristic
polynomial over \(K(z)\).
Relabel so \(w_0<w_1<w_2<w_3\).
Write its polynomial as
\(X^4+c_1X^3+c_2X^2+c_3X+c_4\), and use the \(z\)-valuation.
The unique lowest terms give
\[
 v(c_1)=w_0,\quad v(c_3)=w_0+w_1+w_2,\quad
 v(c_4)=w_0+w_1+w_2+w_3.
\]
If the \(01\) principal \(2\)-minor of \(x\) is nonzero,
\(v(c_2)=w_0+w_1\). The Newton polygon has four
length-one edges with distinct slopes, giving four distinct
roots.
Otherwise every possible term of \(c_2\) has valuation at least
\(w_0+w_2\), which is greater than
\((2w_0+w_1+w_2)/2\).
There are the two outer length-one edges and one middle edge
of length two.
After adjoining \(z^{1/2}\) if necessary, its residual
polynomial is \(Y^2+c\), with \(c\ne0\).
It is separable because \(2\ne0\), so Hensel's lemma again
gives four distinct roots.
This proves the first case without a quartic discriminant
expansion.

In the remaining case the coset interchanges the two
irreducible \(2\)-blocks.
Its matrices have the form
\(\begin{psmallmatrix}0&B\\C&0\end{psmallmatrix}\).
Irreducibility of each block and the same nonzero-function
argument allow all entries of \(B,C\) to be chosen nonzero;
both blocks are invertible.
Order the two cocharacter weights in each block as
\(u_0<u_1\) and \(v_0<v_1\).
The characteristic polynomial after multiplying by \(D(z)\)
has the form
\[
 X^4-A(z)X^2+\Delta(z).
\]
Here the unique lowest term of \(A(z)\) is a nonzero product
of two block entries times \(z^{u_0+v_0}\), and
\[
 v(\Delta)=u_0+u_1+v_0+v_1>2v(A).
\]
The quadratic in \(X^2\) has two distinct nonzero roots of
different valuations. Taking their square roots gives four
distinct roots since \(2\ne0\).

Thus the separable locus is a nonempty Zariski open in every
coset of \(\mathcal H\).
The Zariski closure of \(r(U)\), for a finite-index \(U\),
contains \(\mathcal H\): finitely many translates of that
closure cover \(\mathcal G\), so its identity component has
the same dimension.
Consequently \(r(\sigma U)\) meets the nonempty separable
open just found. This proves the lemma.
\end{proof}

\begin{srproposition}[The actual rank-four curve conditions]
\label{prop:sr-four-curve-genericity}
Let an actual ordinary rank-four curve over \(k[[T]]\) have
absolutely irreducible generic representation.
Suppose its four ordinary characters are pairwise distinct
on each \(U_v\), its residual packet satisfies condition
\(\mathrm{(C)}\) of Lemma~\ref{lem:sr-rank-four-no-self-twist},
and one residual character ratio has order greater than \(4\).
Then it has:
\[
 \begin{gathered}
 \kappa_{v,i}\ne\kappa_{v,j},\qquad
 \kappa_{v,i}\ne\epsilon_\ell\kappa_{v,j}\quad(i\ne j),\\
 r|_{\ker\epsilon_\ell}\text{ absolutely irreducible},\\
 \text{a regular semisimple element in every open coset of }
                         \ker\epsilon_\ell .
 \end{gathered}
\]
In particular these give the rank-four local genericity,
cyclotomic irreducibility, and regular-Frobenius-coset inputs
to the support patching theorem.
\end{srproposition}
\begin{proof}
Both local inequalities follow by restriction to \(U_v\),
where \(\epsilon_\ell=1\).
Lemma~\ref{lem:sr-four-free-local-element} gives an element
\(u\) in one such subgroup whose ordinary eigenvalue ratios
all have infinite order; its characteristic polynomial and
every power's characteristic polynomial are separable.

The proof of Lemma~\ref{lem:sr-rank-four-no-self-twist}
works equally over the algebraic closure of \(k((T))\):
normal-subgroup semisimplicity uses the socle argument, and
the quotient by the cyclotomic kernel is still procyclic.
A nontrivial isotypic orbit has even length dividing \(4\),
hence gives a quadratic self-twist factoring through the
mod-\(\ell\) cyclotomic character. Its values \(1,-1\)
remain distinct after reduction.
Condition \(\mathrm{(C)}\) excludes it.
A single isotypic component with multiplicity greater than
one is excluded by the tensor-factor eigenline argument
and closedness of subspace stabilizers.
This proves cyclotomic irreducibility in characteristic
\(\ell\), without assuming it from the characteristic-zero
case.

Apply Lemma~\ref{lem:sr-four-noncommutative} with
\(Z=\ker\epsilon_\ell\).
It shows that no open normal subgroup of \(Z\) has
commutative image.
The regular-coset lemma now applies to \(r|_Z\).
Polynomial zero sets are closed because matrix entries lie
in the valued power-series coefficient field; a finite
coefficient extension containing a polynomial's coefficients
suffices for this assertion.
This gives the stated regular-coset condition.
\end{proof}

% Frozen-source correspondence:
% WeakShimuraTheory/CompactIntegralFields/H4.lean;
% WeakShimuraTheory/StabilizerIndexBounds/
% ExistsMemPowerSeriesCharRatiosNotIsOfFinOrder.lean, NonComm.lean;
% AlgebraTools/MonomialRepresentations/
% MonomialOfComm.lean, MonomialRed.lean, PermChar.lean, RegCosetFour.lean,
% DiscSwap.lean, RelWeight.lean, TorusTransfer.lean.
% The Newton-polygon argument replaces the source's explicit
% quartic discriminant algebra; it proves the same separability claim.

\subsection{The Case II Hilbert seed and its induction}
\label{subsec:sr-case-two-seed}
We separate the mixed Hilbert seed from its elementary induction.
The latter is ordinary induction from a real quadratic field; the
source name for its matrix change of basis does not denote a
symmetric-square functor.

\begin{srlemma}[Mixed ordinary Hilbert seed]
\label{lem:sr-mixed-hilbert-seed}
Let \(k_0\) be a real quadratic field in which the odd prime
\(\ell\) splits. Let \(\xi:G_{k_0}\to\overline{\mathbb Q}_\ell^\times\)
be finite and totally odd, unramified at \(\ell\), of order prime
to \(\ell\), with \(\bar\xi^2\ne1\).
Suppose its conductor is supported above \(2\), with a ramified
prime of norm \(2\), and its Frobenius value at each prime above
\(\ell\) is nontrivial modulo \(\ell\).
Write these two primes as \(v_+,v_-\).

Using the classical Hilbert Eisenstein, integral q-expansion and
bounded-denominator, ordinary control, and local--global
compatibility theorems at this tame conductor,
there are \(k,t\ge2\), both divisible by \(\ell-1\), and a
cuspidal Hilbert modular representation with associated
\(r:G_{k_0}\to\mathrm{GL}_2(E)\), for a finite
\(E/\mathbb Q_\ell\), such that
\[
 \det r=\epsilon_\ell^k\xi,\qquad
 \bar r^{\,\mathrm{ss}}=1\oplus\bar\xi.
\]
It is irreducible on every open subgroup. Its ordinary weights
and residual ordering are
\[
 \begin{array}{c|c|c}
 &\text{weights}&\text{ordered residual characters}\\\hline
 v_+&(k+t,-t)&(\bar\xi,1)\\
 v_-&(k,0)&(1,\bar\xi).
 \end{array}
\]
Away from \(\ell\), its restriction to inertia is conjugate to
\(1\oplus\xi\).
The same construction permits arbitrary sufficiently large \(k,t\).
\end{srlemma}
\begin{proof}
Let \(\omega\) be the finite Hecke character corresponding to
\(\xi\), with Artin reciprocity convention chosen accordingly.
The parallel weight-one Eisenstein series \(E_1(1,\omega)\)
has normalized nonconstant ideal coefficients
\[
 a(\mathfrak m)=\sum_{\mathfrak a\mid\mathfrak m}
                         \omega(\mathfrak a),\qquad a(1)=1.
\]
In particular these coefficients are integral.
At each prime above \(\ell\), the two old-space ordinary roots
are \(1,\omega(v)\).
Stabilize with the second root at \(v_+\) and the first at
\(v_-\). Denote the resulting form by \(E\).
Equivalently apply \(1-V_{v_+}\) and
\(1-\omega(v_-)V_{v_-}\), with the usual ideal-shift
normalization of \(V\).
Its nonconstant coefficients remain integral and its first is
\(1\).

Choose a parallel weight \(k\), divisible by \(\ell-1\), and
a level-one Hasse lift \(H\) whose expansion at every cusp is
\(1\) to sufficiently high \(\ell\)-adic precision.
Such lifts exist with arbitrarily large weight and precision.
The finitely many constant coefficients of \(E\) have bounded
denominators. Choose the precision larger than those denominators.
Then the nonconstant coefficients of \(EH-E\) are in the
coefficient maximal ideal, and \(EH\) has weight \(k+1\).
No assertion that the weight-one constant coefficients themselves
are integral is needed.

The Hilbert weight-one Fourier formula is classical; a
real-quadratic explicit statement is
\href{https://webpages.leeu.edu/rmoy/files/2018/08/noncm.pdf}
{Moy--Specter, Proposition 2.1}.
The level-one Hasse lift and arbitrary precision follow from
\href{https://arxiv.org/pdf/1911.11196}
{Deo--Dimitrov--Wiese, Lemmas 2.2--2.3}.

For primes \(\mathfrak p\nmid\ell\), the Hecke coefficient
formula changes on passing from weight \(1\) to weight \(k+1\)
only by the factor \(\mathrm N\mathfrak p^k\), which is \(1\)
modulo the maximal ideal. At the Iwahori primes above \(\ell\),
the \(U\)-formula is the same coefficient shift.
Thus \(EH\) is a simultaneous eigenvector modulo that ideal
with the chosen two ordinary roots.
Use the Hecke-stable lattice defined by its nonconstant Fourier
coefficients in the finite-dimensional holomorphic space.
It is a genuine lattice: finitely many coefficient evaluations
are injective, so the integrality conditions give a bounded
module, and the integral q-expansion/bounded-denominator theorem for
classical Hilbert forms supplies a basis with bounded coefficient
denominators, hence a full sublattice. This integral theorem is
an explicit published input here, not a consequence of finite
dimensionality alone.
The commuting finite Hecke algebra and the elementary
Deligne--Serre lifting lemma give a characteristic-zero
holomorphic eigenform \(g\) with the same residual eigensystem.
Its two \(U\)-eigenvalues are units, with residues
\(\bar\omega(v_+)\) and \(1\).

We check cuspidality, which is essential here.
If \(g\) were Eisenstein, its unramified eigenvalues would come
from two global finite characters \(\eta_1,\eta_2\), with
\(\eta_1\eta_2=\omega\), one carrying the factor
\(\mathrm N\mathfrak p^k\).
Chebotarev and characteristic polynomials identify their
residual sum with \(1\oplus\bar\omega\).
Hence the globally unordered residual pair is
\(\{1,\bar\omega\}\).
Since \(k>0\), precisely one of its two local roots at every
\(\ell\)-prime is a unit: the character without the positive
\(\ell^k\) factor. It is the same globally chosen character at
both primes.
But at \(v_+\) it must reduce to \(\bar\omega(v_+)\ne1\),
whereas at \(v_-\) it must reduce to \(1\ne\bar\omega(v_-)\).
Neither global choice is possible. Thus \(g\) is cuspidal.

Apply nearly ordinary control and independence of algebraic
weight, in the component of \(g\), with fixed central character.
Choose the arithmetic weight giving the displayed two weight
pairs, with \(t\) divisible by \(\ell-1\).
Both pairs have sum \(k\), so this holds the determinant fixed.
Their gaps are positive and can be chosen arbitrarily large,
so ordinary classicality applies.
For this application the ordinary module is an actual
definite-quaternion module, not a formal instance of an
axiomatic control statement.
Take the quaternion algebra over the real quadratic field
\(k_0\) ramified at both real places and at no finite place.
It exists because the number of ramified places is even.
At finite places identify its maximal compact subgroups
with those of \(\mathrm{GL}_2(k_{0,v})\), using the same
tame level \(\Gamma_1(\mathfrak n)\) and the same Iwahori
and deeper ordinary levels above \(\ell\).
The coefficient module consists of functions on its finite
double-coset set with the fixed central character.
The local matrix translations, lower-unipotent averaging
and ordinary projector are the rank-two instance of the
ordinary function-module construction.
That construction verifies the projectivity and transition
axioms required for control; the finite central-character
projector is exact since its relevant diamond orders are
powers of \(2\).
The algebraic coefficient at an infinite place has highest
weight difference equal to the prescribed positive ordinary
gap minus one.
Evaluation on its highest-weight line, with the normalized
\(U(y)\) operators, gives weight independence.
At classical regular weights, Jacquet--Langlands identifies
the resulting eigensystems with the cuspidal Hilbert
eigensystems at precisely these finite levels.
Thus this actual module verifies the hypotheses of the
ordinary control theorem.
Its finite torsion-free weight algebra is faithful on its
nonzero localized ordinary space; lying over supplies an
arithmetic point above the chosen weight.
The normalized \(U\)-characters retain the two residual
orientations because \(\bar\xi(\operatorname{Frob}_{v_\pm})\ne1\).
The control and weight normalization are the actual
definite-quaternion instances of
\href{https://www.math.ucla.edu/~hida/Lum07.pdf}
{Hida, Theorems 1.2 and 2.4}.
Theorem 1.2 requires its module axioms; the preceding
function-module construction verifies them here.
The operator \(U(y)\) is divided by its lowest-weight factor
before applying control.

The classical Hilbert Galois representation theorem gives an
irreducible ordinary \(r\) with the stated determinant and
residual semisimplification.
The tame level throughout is exactly
\(\mathfrak n=\operatorname{cond}(\xi)\):
the Eisenstein series starts there, stabilization adds only
the two primes above \(\ell\), the Hasse lift has level one,
and ordinary weight variation keeps this tame level.
If one realizes the spaces at a neat level, it is
\(\Gamma_1(\mathfrak n)\) already within this conductor;
no extra tame prime is introduced.

Here is the local conductor check.
At a prime where \(\xi\) is ramified, write \(V\) for the
two-dimensional Weil--Deligne representation.
Its conductor is at most the level exponent, which is the
conductor exponent of \(\det V=\xi\).
If \(V^I=0\), then
\(a(\det V)\le a(V)-1\):
in the Swan integral the codimension for the determinant is
at most that for \(V\) at every ramification subgroup, while
the tame codimension is strictly smaller.
This is impossible. Consequently \(\dim V^I=1\), and the
inertial representation is \(1\oplus\xi\).
Monodromy is zero, since it commutes with inertia and no
nonzero nilpotent preserves these two distinct one-dimensional
characters. At the other tame primes the level exponent is
zero and the representation is unramified.
Local--global compatibility proves the asserted exact inertia
type.

We finish by proving the stronger open-subgroup irreducibility
rather than adding it as an input.
If \(r\) becomes reducible on a normal open subgroup, Clifford
theory gives either two distinct one-dimensional constituents
or a scalar restriction.
In the former case the two constituents are interchanged,
giving a nontrivial quadratic self-twist \(\eta\).
Since \(\ell\) is odd, reduction preserves its nontriviality.
Then
\[
 \bar\eta(1\oplus\bar\xi)\simeq1\oplus\bar\xi
\]
forces either \(\bar\eta=1\), or
\(\bar\eta=\bar\xi\) and \(\bar\xi^2=1\), both impossible.
In the scalar case the projective image is finite.
The ordinary inertia ratio at \(v_-\) is
\(\epsilon_\ell^k\), of infinite order, excluding this case.
An arbitrary open subgroup contains a normal open subgroup,
so every open restriction is irreducible.
\end{proof}

\paragraph{Published input normalization.}
The preceding lemma uses the classical Hilbert ordinary theory
at the original tame conductor, not an automorphy lifting
theorem for arbitrary reducible residual representations.
The weight-one construction and the two opposite unit roots
produce the initial cuspidal eigensystem explicitly.
The Hasse lift has level one, so it adds no tame prime.
At small \(\ell\), one uses the neat \(2\)-power level already
present in the principal Case II datum and then its finite
character projector; the relevant diamond-group orders are
powers of \(2\), hence units at odd \(\ell\).
The local type assertion must be read in the same
local--global convention as the Hilbert Galois representation.

\begin{srproposition}[The induced Case II seed]
\label{prop:sr-induced-case-two-seed}
Consider an actual Case II tower \(F,L=FK\), and its real
quadratic field \(k_0\), split at \(2\) and \(\ell\).
Let \(\varphi\) be the quadratic character of \(k_0/\mathbb Q\).
On \(G_F\), its residual value is the quadratic root character
in the four-character packet.
Assume the tower root conditions: this character is nontrivial
on \(G_L\), and the packet has no nontrivial quadratic
self-twist on \(G_L\) factoring through mod-\(\ell\)
cyclotomic image.
Let the finite character \(\xi\) be the restriction of the
auxiliary rational character \(\xi:G_{\mathbb Q}\to E^\times\)
of the Case II datum. Let \(r\) be its mixed Hilbert seed from
Lemma~\ref{lem:sr-mixed-hilbert-seed}, with the two ordinary
patterns interchanged by the nontrivial automorphism of
\(k_0/\mathbb Q\).

Then, after a finite coefficient extension, ordinary induction,
a specified symplectic change of basis, and a unique character
twist congruent to \(1\) give the linked rank-four Case II seed:
it is irreducible on \(G_L\cap\ker\epsilon_\ell\), has the
prescribed multiplier and ramification conditions, and its
residual point is in the partially split good stratum.
Its ordinary squared-weight gaps are at least \(4\).
After an explicitly removable scalar twist it is attached to a
regular algebraic cuspidal representation over \(L\).
\end{srproposition}
\begin{proof}
Put \(H=G_{k_0}\), let \(t_0\) be the tower's distinguished
element outside \(H\), and set
\[
 \mu=\epsilon_\ell^k\xi\varphi
       \quad\text{on }G_{\mathbb Q}.
\]
The determinant of \(r\) is \(\mu|_H\).

First choose a Ribet lattice on \(H\cap G_F\).
This restriction is irreducible, whereas its residual
semisimplification is \(1\oplus\bar\xi\).
The values of these characters at \(t_0^2\) are distinct.
Thus the stable-lattice fixed set in the Bruhat--Tits tree lies
in the apartment determined by the two eigenlines of
\(r(t_0^2)\).
It is a bounded interval: an unbounded interval would give an
invariant line over the fraction field.
At its endpoint with upper residual orientation, reduction is
an extension of \(\bar\xi\) by \(1\) that is nonsplit on
\(H\cap G_F\); otherwise the interval could be extended.
In the corresponding eigenbasis,
\[
 r(t_0^2)=\operatorname{diag}(\alpha^2,\beta^2),\qquad
 \alpha\beta=\mu(t_0),
\]
and all matrices on \(H\cap G_F\) are integral, reducing to
\(\begin{psmallmatrix}1&b(h)\\0&\bar\xi(h)\end{psmallmatrix}\),
with \(b(h)\ne0\) for some \(h\).
Choose the square root \(\alpha\) in the prescribed residue
class \(\pm1\); Hensel's lemma applies since \(\ell\ne2\).
Then put \(\beta=\mu(t_0)/\alpha\).
The residual packet identity fixes its required residue too.

In the coset basis, \(I=\operatorname{Ind}_H^{G_{\mathbb Q}}r\)
has matrices
\[
 I(h)=
 \begin{pmatrix}r(h)&0\\0&r(t_0^{-1}ht_0)\end{pmatrix},
 \qquad
 I(t_0)=
 \begin{pmatrix}0&r(t_0^2)\\1&0\end{pmatrix}.
\]
Write \(J_2=\begin{psmallmatrix}0&1\\-1&0\end{psmallmatrix}\).
The alternating form
\(\operatorname{diag}(J_2,\mu(t_0)J_2)\) has multiplier
\(\mu\) under both displayed kinds of matrices.
For the second matrix this follows from
\(\det r(t_0^2)=\mu(t_0)^2\).
Define
\[
 P=
 \begin{pmatrix}
 0&\alpha&0&\alpha\\
 \beta&0&-\beta&0\\
 0&-1&0&1\\
 1&0&1&0
 \end{pmatrix}.
\]
It is invertible because \(2\alpha\beta\ne0\), and
\[
 P^t\operatorname{diag}(J_2,\alpha\beta J_2)P
       =-2\alpha\beta J_4,\qquad
 J_4=\begin{pmatrix}
 0&0&0&1\\0&0&1&0\\0&-1&0&0\\-1&0&0&0
 \end{pmatrix}.
\]
Consequently \(R=P^{-1}IP\) is a symplectic similitude of
multiplier \(\mu\), and
\[
 R(t_0)=\operatorname{diag}(\beta,-\alpha,-\beta,\alpha).
\]
These equations specify both the normalization and the pairing.

We prove the required cyclotomic irreducibility.
Put \(V=G_L\) and \(U=V\cap\ker\epsilon_\ell\).
If \(\varphi|_U=1\), then \(\varphi|_V\) factors through the
closed cyclotomic image.
Its quadratic character kills the pro-\(\ell\) subgroup, so
it factors through the mod-\(\ell\) image.
But multiplication by \(\varphi\) permutes the residual packet
\(\{1,\bar\xi,\varphi,\bar\xi\varphi\}\).
The no-self-twist condition would force \(\varphi|_V=1\),
contrary to the other root condition.
Thus \(U\not\subset H\).

The restriction \(r|_{H\cap U}\) is irreducible:
apply the strong-open irreducibility and procyclic-restriction
lemma to the open group \(H\cap V\) and its cyclotomic kernel.
Suppose its conjugate by \(t_0\) were isomorphic to it, with
intertwiner \(M\).
For every \(h\in H\cap V\), normality of \(H\cap U\) and
Schur's lemma make
\[
 r(h)^{-1}M^{-1}r(t_0^{-1}ht_0)M
\]
a scalar. The determinants of the two representations agree,
so that scalar has square \(1\).
It follows that \(r(h^2)\) and \(r(t_0^{-1}h^2t_0)\) have
the same characteristic polynomial.

Choose \(h\) in inertia above \(\ell\) and in \(H\cap V\)
with \(\epsilon_\ell(h)\) of infinite order.
At the two conjugate primes, the eigenvalue pairs of these
squares are respectively
\[
 \{\epsilon_\ell(h)^{2(k+t)},\epsilon_\ell(h)^{-2t}\},
 \qquad
 \{\epsilon_\ell(h)^{2k},1\}.
\]
They differ: \(1\) belongs only to the second pair.
Thus the two irreducible constituents on \(H\cap U\) are
nonisomorphic.
The elementary index-two induction criterion now gives
\(R|_U\) irreducible.

Here is the residual good-point calculation.
For \(h\in H\cap G_F\), the two blocks reduce to
\[
 A=\begin{pmatrix}1&b\\0&x\end{pmatrix},\qquad
 B=\begin{pmatrix}1&b'\\0&x\end{pmatrix},
 \qquad x=\bar\xi(h).
\]
Put \(\gamma=\bar\beta/\bar\alpha\).
Direct multiplication by \(P^{-1}\) and \(P\) gives
\[
 \bar R(h)=
 \begin{pmatrix}
 x&0&0&0\\
 (\gamma b-b')/2&1&-(\gamma b+b')/2&0\\
 0&0&x&0\\
 (\gamma b+b')/2&0&-(\gamma b-b')/2&1
 \end{pmatrix}.
\]
Its off-diagonal support is
\(\{-\alpha_{\rm root},\beta_{\rm root},-\theta_{\rm root}\}\)
in the root notation of the preceding subsection.
At a point with \(b\ne0\), either \(\gamma b-b'\ne0\),
or \(\gamma b+b'=2\gamma b\ne0\).
Hence either the negative-\(\alpha_{\rm root}\) entry is
nonzero, or both the \(\beta_{\rm root}\) and
negative-\(\theta_{\rm root}\) entries are nonzero.
The outside coset differs by the normalized diagonal
\(R(t_0)\), so the same support and the prescribed diagonal
packet hold on all of \(G_F\).

We also check the local flag rather than inferring it merely
from induction.
At the prime with weights \((k+t,-t)\), the high line in
the first block has residual eigencharacter \(\bar\xi\);
write a primitive vector on it as \((u,1)\).
The high line in the other block, with weights \((k,0)\),
has residual character \(1\), so it reduces to its first
coordinate line.
In the mixed basis these two vectors have coordinates
\[
 (1/(2\beta),u/(2\alpha),-1/(2\beta),u/(2\alpha)),
 \qquad (0,-1/2,0,1/2)
\]
modulo the maximal ideal, up to units.
Their first two successive flag minors are units.
At the conjugate place the roles are reversed, giving
\[
 (1/2,-u/2,1/2,u/2),\qquad
 (0,1/(2\alpha),0,1/(2\alpha)).
\]
Again the two minors are units.
The symplectic flags are therefore in the integral lower open
cell. Their unique lower unipotent representatives satisfy
\[
 n_{v,20}\equiv-1\text{ or }1,\qquad n_{v,21}\equiv0.
\]
The third step is the symplectic perpendicular of the first,
so these two minors suffice for the complete symplectic flag.
Together with the displayed residual matrix this is exactly
the partially split good stratum.

Let \(\vartheta\) be the required tower multiplier.
The residual packet gives
\(\vartheta/\mu\equiv1\) on \(G_F\).
Since squaring is an automorphism of \(1+\mathfrak m_E\)
for odd \(\ell\), there is a unique continuous
\(\delta:G_F\to1+\mathfrak m_E\) with
\(\delta^2=\vartheta/\mu\), after enlarging \(E\) to contain
the finitely many character values.
Set \(Q=\delta R|_{G_F}\).
It has multiplier \(\vartheta\), the same residual point,
and the same ordinary flags.
Its local torsion condition holds because cyclotomic images
have no \(\ell\)-power torsion, the other local diagonal
characters are unramified, and \(\delta\equiv1\).

For the global ramification condition, take an allowed inertia
element away from \(\ell\) with trivial prescribed residual
representation. Prime-to-\(\ell\) reduction gives
\(\xi=\psi=\varphi=1\) there.
Both rank-two blocks are \(1\) by their actual inertia type.
Thus \(R=1\), \(\vartheta/\mu=1\), and the unique square root
near \(1\) is \(\delta=1\).
So \(Q\) kills the required normal subgroup.

Finally put \(M=Lk_0\).
Solvable base change gives the cuspidal rank-two representation
over \(M\); its irreducible Galois representation rules out
cuspidal splitting.
The unequal ordinary weight pairs also rule out invariance
under \(\operatorname{Gal}(M/L)\).
Quadratic automorphic induction is therefore cuspidal over \(L\).
The union of weights is
\[
 m=(k+t,k,0,-t),
\]
with each consecutive gap at least \(2\), so this induction
is regular algebraic after the standard integral norm twist.
For example, under
\(r_\ell(\pi)=\operatorname{rec}_\ell
(\pi\otimes|\det|^{(1-n)/2})\),
one takes
\(\operatorname{AI}_{M/L}(\pi_M)\otimes|\det|\);
its attached representation is exactly \(R|_{G_L}\).
This checks the rank \(2\) to rank \(4\) normalization.

In the actual principal datum
\(\vartheta=\epsilon_\ell^{a+1}\omega^{-1}\xi\psi\).
Hence
\[
 \delta^2=\epsilon_\ell^{b}\,\omega^{-1}\psi\varphi^{-1},
 \qquad b=a+1-k,
\]
whose second factor has finite order prime to \(\ell\).
The squared weights of \(Q\) are \(2m_i+b\);
their consecutive gaps are at least \(4\).
Twisting \(Q|_{G_L}\) by \(\delta^{-1}\) removes the correction
and gives the actual attached automorphic \(R|_{G_L}\).
Moreover
\((\delta^{-2}\epsilon_\ell^b)^N=1\) for a positive \(N\),
and \(-b+(2m_0+b)\) is even.
These are exactly the finite-order and parity conditions
required of the seed's removable scalar twist.

All matrices and characters used are defined over a finite
coefficient field: \(r\) is, \(\alpha,\beta\) require only a
finite extension, the unique square root near \(1\) remains
there, and the local flag vectors are obtained by solving
finite linear equations with unit flag minors.
Thus the integral seed has the required finite coefficient
DVR, rather than only a completed algebraic-closure lattice.
\end{proof}

% Frozen-source correspondence (alpoge-source-7cbf7a0c):
% ModularityArgument/ThresholdArguments/DenseThresholdCases/
% MPropMHol.lean, THolLift.lean, SVRPropMV.lean, BirthMix.lean,
% MixSeedOfHolWC.lean, MixSeedOf.lean, SeedIIZ.lean, BirthSeedII.lean;
% BuildingsResidueChars/SymmetricSquareMixedExists.lean,
% ExistsRibetLattice.lean, GaloisCharacterResidualDecomposition.lean,
% ResidualCharactersPsiNe.lean, MixSignTwist.lean,
% InducedIrreducible.lean, InducedRepresentationConstituentVanishing.lean;
% Foundations/GaloisRepresentations/AdmissibleResidues/IndLoc0.lean.

\subsection{Local vanishing of the modular ideal at the actual seed}
\label{subsec:sr-seed-local-zero}

\begin{srlemma}[A faithful finite quotient and zero relative tangent]
\label{lem:sr-faithful-quotient-tangent}
Let \(W\) be a regular weight algebra in characteristic zero,
\(D\) a Noetherian deformation algebra, and \(J\subset D\).
Suppose \(D/J\) is finite and torsion-free over \(W\).
Let \(x:D\to E\), with \(E\) a finite coefficient field,
kill \(J\).
If every first-order deformation of \(x\) fixing its weight
is constant, then
\[
 J D_{\ker x}=0.
\]
In particular every prime \(\mathfrak q\subset\ker x\)
contains \(J\).
\end{srlemma}
\begin{proof}
Localize at the weight prime and at \(\ker x\), and complete.
After a finite separable residue-field extension, lifted
uniquely by finite etale extension of the complete weight
ring, the two residue fields agree.
Call the resulting complete weight ring \(W'\) and the
completed deformation ring \(S\).
The zero relative tangent condition says
\(\mathfrak m_S=\mathfrak m_{W'}S+\mathfrak m_S^2\).
Formal Nakayama, or successive approximation modulo powers of
the maximal ideal, makes \(W'\to S\) surjective.

The corresponding local factor \(T\) of the completion of
\(D/J\) is a nonzero finite torsion-free \(W'\)-module.
It is nonzero because \(x\) kills \(J\).
Torsion-freeness persists under completion and the finite
etale extension, and under passage to the local direct factor.
Since \(W'\) is a domain, \(W'\to T\) is injective.
But it factors as the surjection \(W'\to S\) followed by
\(S\to T\). The first map is therefore also injective, and
the second has zero kernel.
Thus the localized completed ideal is zero.
Faithful flatness of completion gives \(J D_{\ker x}=0\).

For \(j\in J\), choose \(s\notin\ker x\) with \(sj=0\).
If \(\mathfrak q\subset\ker x\), then \(s\notin\mathfrak q\),
so primality forces \(j\in\mathfrak q\).
\end{proof}

\begin{srproposition}[Seed local vanishing]
\label{prop:sr-seed-local-zero}
Let \(J_0\) be the actual unitary ordinary Hecke ideal for a
linked tower over \(L=FK\).
Use its finite torsion-free weight quotient and its classical
arithmetic-point property.
Let \(x\) be the arithmetic point obtained from the induced
Case II seed, after the harmless scalar twist making its
weights integral in the unitary normalization.
Then every prime below \(\ker x\) contains \(J_0\).
This remains true after transporting \(J_0\) to the common
coefficient algebra and changing that harmless scalar twist.
\end{srproposition}
\begin{proof}
The seed is an actual ordinary automorphic representation
with regular weights, the prescribed tame inertial type,
and unramified representation at the extra link place.
The classical arithmetic-point property therefore puts it
in the Hecke quotient: \(J_0\subset\ker x\).
This uses its actual automorphy and level, not a proposed
lifting theorem.

Consider a first-order deformation of \(x\) with weight fixed.
At an absolutely irreducible point, the determinant
deformation is represented by a matrix representation over
\(E[\tau]/\tau^2\), unique up to strict conjugation.
One elementary construction lifts matrix units in the
Cayley--Hamilton algebra: their reductions generate \(M_4(E)\)
by absolute irreducibility, and their multiplication errors
are corrected across a square-zero ideal.
The polarization relation lifts with them.
Write the resulting representation as
\[
 r'(g)=(1+\tau z(g))r(g).
\]
Its multiplication law is the adjoint cocycle equation.
The lifted conjugate-self-dual pairing extends this cocycle
to the descended polarized adjoint on \(G_F\).

The restriction at every \(\ell\)-place has an actual lifted
ordinary flag. To check this for an ordinary determinant
deformation, choose an inertia element after killing the
finite characters, with cyclotomic value of infinite order.
The distinct ordinary weights make its four eigenvalues
distinct. Their differences are units over
\(E[\tau]/\tau^2\), so its spectral idempotents lift.
The ordinary Cayley--Hamilton relations kill the lower
matrix entries between these ordered idempotents and
identify each diagonal character; they give the required
free flag. This is why an ordinary determinant relation
suffices here.
Fixing weight makes its diagonal inertial characters constant.
The consecutive seed gaps are at least \(4\) for the squared
weights, or at least \(2\) after the integral unitary
normalization. In particular no upper character ratio is
the cyclotomic character.
The positive-ordinary-extension lemma therefore makes the
first-order local representation crystalline after killing
its finite inertial characters.
Finite-index restriction is injective on \(H^1\) in
characteristic zero by averaging.
Thus its local cocycle lies in the Bloch--Kato finite
subspace before that extension too.

Away from \(\ell\), the deformation relations kill the
inertia subgroup on which the prescribed residual characters
are all trivial.
The seed's actual inertia representation is a direct sum of
those finite characters.
After this subgroup is killed, averaging over the remaining
finite quotient makes the inertial cocycle a coboundary.
At the additional unramified link place the same statement
is immediate. Hence the global cocycle belongs to
\[
 H^1_f\bigl(F,\operatorname{ad}_{\rm pol}(r|_{G_L})\bigr),
\]
the adjoint descended by its actual polarization.

All hypotheses of the published adjoint theorem are now
present: the seed over \(L\) is regular algebraic cuspidal,
its cyclotomic-kernel restriction is irreducible, and the
polarization just used is the one from its automorphic
conjugate self-duality.
Thorne's Theorem 6.1 gives vanishing of this descended
adjoint group over the totally real field \(F=L^+\).
It does not assert full adjoint vanishing over \(L\).
If the seed is first written with the correction character
\(\delta\), remove it using \(\delta^{-1}\), which extends
to \(G_F\); this leaves the adjoint unchanged and matches
the actual automorphic representation and polarization.
We conclude that \(z\) is a coboundary, so \(r'\) is strictly
conjugate to \(r\).

The ordered ordinary characters also have zero derivatives.
They cannot vary inside a constant representation: locally
they are pairwise distinct as characters, so the successive
invariant lines with the prescribed ordering are unique.
Equivalently, the infinitesimal flag deformation maps between
distinct character constituents and has no invariants.
Thus all characteristic-polynomial coefficients and all
ordered ordinary characters are constant.
They generate the universal ordinary determinant algebra;
the first-order map \(x\) is constant.

Apply Lemma~\ref{lem:sr-faithful-quotient-tangent} to the
actual finite torsion-free Hecke quotient.
This gives the local vanishing assertion.

For transport to the common space, a prime below the
seed's common-space point pulls back to a prime below \(x\),
hence contains \(J_0\).
Its image ideal is therefore contained in the original prime.
For a scalar change of normalization, twist the universal
determinant and ordinary characters by the corresponding
finite/algebraic character.
The same twist is an automorphism of the ordinary Hecke
module, with inverse the inverse-character twist.
It preserves its annihilator. Arithmetic-point density and
adic closedness identify the transported closed Hecke
ideals, proving the final assertion.
\end{proof}

\paragraph{Assembly with the two-link construction.}
The actual modular ideal used for Case II is the sum of the
images of two such unitary Hecke ideals, formed from links
with the same scalar normalization and two different
auxiliary unramified places.
Ordinary finite-weight control gives
\eqref{eq:sr-pm-finite}.
The rank-independent Taylor--Wiles support theorem applies
to its characteristic-\(\ell\) curves:
Proposition~\ref{prop:sr-four-curve-genericity} provides the
rank-four genericity, cyclotomic irreducibility, and regular
coset hypotheses; the local link condition gives the
additional unramified-place condition.
Proposition~\ref{prop:sr-seed-local-zero} gives the seed
property for each ideal, hence for their sum.
Arithmetic ordinary control gives purity and automorphy away
from the corresponding link place.
The two different link places remove the omitted-place
qualification, and the usual solvable descent gives the
automorphy statement over every allowed imaginary quadratic
restriction where it remains irreducible.
These are the five properties of the actual modular ideal
used in the Case II propagation, with every rank-four
specific input supplied above.
The remaining shared input in this assembly is the
rank-independent support patching theorem and its actual
unitary Hecke-module realization.

% Frozen-source map:
% ModularityArgument/ThresholdArguments/DenseThresholdCases/
% EigenbasisRankFour.lean, BirthPm.lean, HidaModularIdealMap.lean,
% PmTwo.lean, PmExTh.lean, PointTheoremBaseObject.lean;
% WeakShimuraTheory/EigenbasisLifts/RegOfTan.lean, Factor.lean,
% LiftRep.lean, RelForm.lean, CocRel.lean, CocExt.lean, FlagDual.lean,
% BkAbEmb.lean, AwAvg.lean, ConjOf.lean, ConstOf.lean;
% PatchingModules/AutomorphicDetection/PfromAut.lean;
% ModularityArgument/ThresholdArguments/DenseThresholdCases/
% THidaU.lean, ExistsHeckeIdeal.lean, UnitaryHeckeData.lean.

\subsection{The arithmetic Hecke quotient on the Case II tower}
\label{subsec:sr-case-two-hecke}
We identify the automorphic algebra in the rank-four support
calculation with the ordinary modules of
Section~\ref{sec:sr-ordinary-control}.
The coefficient lattice, local types and operator
normalizations must be the same on both sides.

Keep an actual Case II tower. Write \(F\) for its totally real
field, \(L=FK\) for its CM field, and \(S\) for the finite
ramification set, including the places above \(\ell\) and the
distinguished auxiliary place \(v_0\).
The tower conditions give
\[
 L/F\ \hbox{unramified at every finite place},\qquad
 v\in S\Longrightarrow v\ \hbox{splits in }L,\qquad
 [F:\mathbb Q]>1 .
 \tag{H1}
\]
The last assertion follows from the deep layer.
The definite rank-four unitary group is quasi-split at every
finite place: its global existence parity is satisfied since
\(4[F:\mathbb Q]/2=2[F:\mathbb Q]\) is even.

We use classical automorphic theorems in the following scope.
Base change of a classical representation on this definite
unitary group is an isobaric sum of discrete
conjugate-self-dual general-linear representations.
This is Labesse's Corollary 5.3 in
\emph{Changement de base CM et s\'eries discr\`etes}; see
\href{https://arxiv.org/pdf/0906.4189}{Guerberoff, Theorem 2.2}.
Its field, degree and finite quasi-splitness hypotheses
are (H1).
For the cuspidal seed's descent with specified local factors,
one can use
\href{https://arxiv.org/pdf/1409.3731}
{Kaletha--M\'inguez--Shin--White, Theorem 1.7.1}
in its proved generic-parameter, pure-inner-twist case
(the qualification immediately following the theorem).
Our definite hermitian group is a pure inner form.
The simple cuspidal parameter has trivial global component
group modulo the center. Thus one may choose its prescribed
split local members, its spherical nonsplit members, and
its compact cohomological archimedean members.
There is no remaining global multiplicity-character condition.
The existence parity was checked above. No general
nongeneric starred case of that theorem is used.
Attached semisimple Galois representations are supplied by
Guerberoff, Theorem 2.3.
For their cuspidal constituents we use the polarization and
full local--global compatibility of
\href{https://annals.math.princeton.edu/wp-content/uploads/annals-v179-n2-p03-p.pdf}
{Barnet--Lamb--Gee--Geraghty--Taylor, Theorem 2.1.1}.
At \(\ell\), potential semistability and compatibility can be
checked after killing the finite inducing characters, where
the local representation is Iwahori-spherical.
The ordinary local flag has the characters determined by
the unit and uniformizer equations below.
These inputs concern already automorphic representations;
none asserts automorphy of an arbitrary deformation.

\subsubsection{The local lattice and the weight normalization}
Let \(\mathcal O\) be a finite coefficient DVR, enlarged to contain
the needed finite character values, and let
\(\bar\chi_0,\ldots,\bar\chi_3\) be the link's residual
characters on \(G_L\).
At a split tame place \(u\), put
\[
                 \tau_u=\bigoplus_{i=0}^3
                       [\bar\chi_i]|_{I_{L_w}},
 \tag{H2}
\]
where brackets denote finite Teichm\"uller lifts.
This is a finite-image inertial type extending to the Weil group.

The local type used here must detect zero monodromy.
The Schneider--Zink type \(\sigma^{\rm crys}(\tau_u)\)
is finite dimensional and smooth on
\(\mathrm{GL}_4(\mathcal O_{L_w})\), and occurs in a smooth
irreducible \(\pi_w\) exactly when
\(\operatorname{rec}(\pi_w)|_I=\tau_u\) and \(N_{\rm WD}=0\).
A precise formulation is
\href{https://math.uchicago.edu/~emerton/pdffiles/bonn-lectures-survey-article.pdf}
{Emerton--Gee, Theorem 9.2.4}.
Use the contragredient in the function coefficient module:
\[
              \tau_0=\bigotimes_u
                              \sigma^{\rm crys}(\tau_u)^\vee .
 \tag{H3}
\]
Invariants of \(\tau_0\otimes\pi_w\) are then
\(\operatorname{Hom}(\sigma^{\rm crys}(\tau_u),\pi_w)\),
which fixes the dual convention.

There is a finite-free stable integral lattice in (H3)
over a finite coefficient extension.
Each compact representation factors through a finite quotient.
Its finitely many matrices and inverses are defined over
a finite extension; the sum of the finite orbit of a full
lattice is stable.
Tensor these lattices over the finitely many places.
No order of a finite quotient is inverted.

After the prescribed finite local torsion is fixed, the
inertial weight algebra is an actual power-series algebra
\[
 \Lambda\simeq
 \mathcal O[[X_{vij}:v\mid\ell,\ 0\le i<4,\ 1\le j\le d_v]],
           \qquad d_v=[L_w:\mathbb Q_\ell].
 \tag{H4}
\]
It represents the free pro-\(\ell\) unit quotients at
the selected \(w\mid v\).
Write \(\kappa_{vi}\) for its universal inertial characters.
Use the source's local reciprocity convention and define
the torus character \(\theta\) by
\[
 \kappa_{vi}(g)\epsilon_\ell(g)^i
                  \theta(e_{vi}(y))=1,\qquad
            \operatorname{Art}_v(g)=y\in\mathcal O_{L_w}^{\times}.
 \tag{H5}
\]
Here \(e_{vi}(y)\) has \(y\) in coordinate \(i\).
Local reciprocity makes (H5) well-defined and continuous.
The function model uses \(\theta^{-1}\).
Thus the factor \(\epsilon_\ell^i\) is part of the
normalization, and is retained.

For the commuting normalized ordinary operators, indexed
by \(0\le j<4\), the same convention gives
\[
       U_{vj}=\prod_{i=0}^{j}
                  \psi_{vi}(g_v)\epsilon_\ell(g_v)^i,
                  \qquad\operatorname{Art}_v(g_v)=\varpi_v .
 \tag{H6}
\]
Away from \(S\), the residual spherical values are specified by
\[
 \begin{split}
 P_w(X)&=\prod_{i=0}^3
          (X-\bar\chi_i(\operatorname{Frob}_w^{-1}))\\
       &=\sum_{j=0}^{4}(-1)^j
                  q_w^{j(j-1)/2}\bar a_{wj}X^{4-j},
                        \qquad\bar a_{w0}=1 .
 \end{split}
 \tag{H7}
\]
Since \(q_w\) is a unit in the residue field, these
\(\bar a_{wj}\) are uniquely determined.
The residual ordinary values are those of (H6).
Choose arbitrary integral lifts; the simultaneous residual
localization depends only on their reductions.

\subsubsection{Separation by classical specializations}
Let \(M\) be the corresponding residual direct summand of
the constructed ordinary function module, with lattice (H3)
and character (H5).
It is finite free over \(\Lambda\), by the ordinary projector
and coefficient-change proofs, also after the direct-summand
projection. If the residual system does not occur, set \(M=0\).
Otherwise let
\[
              \mathbb T\subseteq\operatorname{End}_\Lambda(M)
 \tag{H8}
\]
be the closed algebra generated by the selected spherical,
ordinary and scalar operators.
It is commutative: this is spherical and ordinary Iwahori
commutativity, and commutativity at disjoint places.
It is finite and torsion free over \(\Lambda\), because
it is a submodule of the finite endomorphism module over
the Noetherian domain \(\Lambda\).

\begin{srlemma}\label{lem:sr-unitary-classical-separation}
The evaluations of \(\mathbb T\) at actual classical ordinary
eigensystems separate its elements.
One may use a single fixed regular algebraic weight,
varying finite-order characters above \(\ell\), and require
the four finite inertial characters at each selected
place to be pairwise distinct.
Consequently \(\mathbb T\) is reduced.
\end{srlemma}
\begin{proof}
In the free coordinates (H4), fix the algebraic part.
Finite-order twists assign roots of unity of \(\ell\)-power
order to a basis of the free pro-\(\ell\) quotient,
after multiplication by the fixed characters for that
algebraic part.
They respect the prescribed finite torsion.
The resulting coordinate values, after subtracting \(1\),
have absolute value less than one.
Each assignment lies in a finite coefficient extension.

These assignments are Zariski dense for integral power series.
In one variable a nonzero element of \(\mathcal O[[X]]\)
has only finitely many zeros in the open unit disc.
Divide out its coefficient-content power of the uniformizer
and apply Weierstrass preparation at the first unit coefficient;
the remaining power-series factor is a unit on that disc.
It cannot vanish at every allowed torsion value.
For several variables, specialize all but one and induct.
The same proof permits finitely many values to be excluded
in each free coordinate.
As \(d_v>0\), it permits all equalities between the four
finite inertial characters to be excluded at each place.

At such a weight, coefficient and classical control identify
\(M\) with actual classical ordinary forms; a specialized
\(\Lambda\)-basis remains linearly independent.
The definite automorphic space is semisimple.
Its unramified operators act by their Satake characters.
At \(\ell\), distinct finite inertial characters exclude
linked inducing characters and monodromy between them.
The ordered Jacquet line has the specified unit and
uniformizer characters.
The commuting normalized ordinary operators act on that
line by one eigencharacter, and act by the same scalars
on repeated automorphic copies.
The selected commuting algebra is therefore a sum of its
eigenspaces on this classical ordinary space.

If \(t\in\mathbb T\) vanishes at every such eigencharacter,
it acts as zero in every such specialization of \(M\).
Each entry of its action matrix in a \(\Lambda\)-basis
vanishes on the dense set just proved, hence is zero.
The faithful inclusion (H8) gives \(t=0\).
A nilpotent element vanishes at every field-valued character,
so \(\mathbb T\) is reduced.
\end{proof}

\subsubsection{Integral interpolation at the prime three}
\begin{srlemma}\label{lem:sr-compact-determinant-interpolation}
Let \(T\) be a compact Hausdorff commutative ring, and
\(x_\alpha:T\to E_\alpha\) a separating family of continuous
maps to characteristic-zero fields.
Let \(\Gamma\) be profinite and
\(r_\alpha:\Gamma\to\mathrm{GL}_n(E_\alpha)\) continuous.
Suppose that on a dense union of conjugacy classes their
characteristic polynomials are specializations of
polynomials in \(T[X]\).
There is then a unique continuous \(n\)-dimensional
determinant on \(T[\Gamma]\) with these specializations.
No integer up to \(n!\) is inverted.
\end{srlemma}
\begin{proof}
The injection \(T\to\prod_\alpha E_\alpha\) is a
homeomorphism onto a closed subring, by compactness and
the Hausdorff property.
For \(g\in\Gamma\), its tuple of characteristic coefficients
is a limit of the corresponding tuples on the dense set.
Conjugation leaves these unchanged.
Thus they belong to \(T\), and their maps \(\Gamma\to T\)
are continuous.

Characteristic polynomials suffice without dividing power
sums by \(1,\ldots,n\).
For finitely many matrices \(A_1,\ldots,A_s\), the integral
Amitsur identity is
\[
 \det(1-\textstyle\sum_i t_iA_i)
       =\prod_{[w]\ {\rm primitive}}\det(1-t_wA_w).
 \tag{H9}
\]
The product is over primitive cyclic words, and \(A_w,t_w\)
are the matrix word and its commuting monomial.
In bounded degree only finitely many factors contribute.
For generic matrices over \(\mathbb Z\), prove this over
\(\mathbb Q\) by formal logarithms and grouping traces
into cyclic words.
Both sides have integral coefficients, so this proves
an identity over \(\mathbb Z\), valid after arbitrary
scalar extension.

Apply (H9) to \(A_i=r_\alpha(g_i)\).
Every factor is a characteristic polynomial of a group
word, whose coefficients already belong to \(T\).
Hence every coefficient on the left belongs to \(T\).
This defines the degree-\(n\) polynomial law on finite
sums in \(T[\Gamma]\).
Multiplicativity, the unit and scalar-extension identities
hold after every separating \(x_\alpha\), and therefore
over \(T\). The same argument proves uniqueness.
The characteristic-coefficient continuity above gives
its profinite continuity.
\end{proof}

Apply this lemma to (H8) and
Lemma~\ref{lem:sr-unitary-classical-separation}.
The finite algebra \(\mathbb T\) over the compact
power-series ring \(\Lambda\) is compact.
Its arithmetic maps are continuous and lie in finite
coefficient extensions, by the automatic-continuity
argument following classical control.
The attached classical semisimple representations give
the \(r_\alpha\).
Chebotarev supplies the dense conjugacy classes and
(H7), with Hecke operators, supplies their polynomials.
This constructs an actual continuous determinant
\(D_{\mathbb T}\).

We record the normalization for arbitrary discrete constituents,
including odd-dimensional ones.
This is the auxiliary-character construction in the proof
of Guerberoff's realization theorem, made explicit in rank four.
Write the discrete base change as
\(\boxplus_i\Pi_i\), of ranks \(n_i\), with
\(\sum_i n_i=4\).
Choose a unitary anticyclotomic Hecke character \(\eta_i\),
unramified at any prescribed finite set including the
places above \(\ell\), whose infinity type is
\[
 \eta_{i,\tau}(z)=(z/|z|)^{\delta_{i,\tau}},
 \quad
 \delta_{i,\tau}=0\ \hbox{if }4-n_i\ \hbox{is even},
 \quad
 \delta_{i,\tau}\in\{1,-1\}\ \hbox{otherwise}.
 \tag{H10a}
\]
The signs at conjugate embeddings are opposite.
The CM Hecke-character construction used in that theorem
supplies these characters with the stated finite
unramifiedness; for even \(n_i\) take \(\eta_i=1\).
Then \(\Pi_i\eta_i\) is cohomological.
The discrete-spectrum classification writes
\[
       \Pi_i\eta_i
             =\boxplus_{j=0}^{b_i-1}
                         \pi_i^0|\det|^j,\qquad n_i=a_i b_i,
\]
where \(\pi_i^0\) is cuspidal of rank \(a_i\) and
\(\pi_i^0|\det|^{(b_i-1)/2}\) is conjugate self-dual.

Choose a second unitary anticyclotomic character \(\xi_i\),
again unramified at the prescribed finite set, with angular
infinity exponent \(0\) if \(a_i\) or \(b_i\) is odd,
and exponent \(1\) or \(-1\) if both are even.
Then
\[
       C_i=\pi_i^0\xi_i|\det|^{(b_i-1)/2}
\]
is cuspidal, cohomological and conjugate self-dual.
The two Hecke characters
\[
 A_i=\xi_i^{-1}|\cdot|^{(a_i-1)(b_i-1)/2},
       \qquad
 B_i=\eta_i^{-1}|\cdot|^{(n_i-4)/2}
 \tag{H10b}
\]
are \emph{algebraic}.
Indeed an angular factor \((z/|z|)^\delta\) has
exponents \(\delta/2,-\delta/2\), while the normalized
complex norm has exponents \(1,1\).
For \(A_i\), the integer
\((a_i-1)(b_i-1)\) is odd exactly when both \(a_i,b_i\)
are even, the case where the angular exponent is odd.
For \(B_i\), \(n_i-4\) is odd exactly when the angular
exponent in (H10a) is odd.
Both pairs of resulting exponents are therefore integers.
Only these algebraic products are converted into
\(\ell\)-adic characters.

Let \(r_{\ell,\iota}(C_i)\) use its rank-\(a_i\)
normalization, and set
\[
 R_i=r_{\ell,\iota}(C_i)\,
            \epsilon_\ell^{\,a_i-n_i}\,
            r_{\ell,\iota}(A_i)\,
            r_{\ell,\iota}(B_i).
 \tag{H10c}
\]
All powers of \(\epsilon_\ell\) in this formula are integers.
At a compatible finite place, the norm exponent in its
parameter is
\[
 \frac{b_i-a_i}{2}+(a_i-n_i)
       +\frac{(a_i-1)(b_i-1)}{2}
       +\frac{n_i-4}{2}=-\frac32 .
\]
The factors \(\xi_i\) cancel and the remaining character
is \(\eta_i^{-1}\).
Thus \(R_i\) has parameter
\(\operatorname{rec}(\pi_i^0\eta_i^{-1}|\det|^{-3/2})\).
Consequently the total attached representation is
\(\bigoplus_i\bigoplus_{j=0}^{b_i-1}R_i\epsilon_\ell^j\).
Putting \(\rho_i'=R_i\epsilon_\ell^{b_i-1}\) and reversing
the index gives exactly the form used below:
\[
             \bigoplus_i\ \bigoplus_{j=0}^{b_i-1}
                                  \rho_i'\epsilon_\ell^{-j}.
 \tag{H10}
\]
This proves the decomposition with actual continuous
finite-coefficient Galois representations.
It does not choose a square root of \(\epsilon_\ell\).
The polarization and Hodge numbers are those of the
definite-unitary realization; the auxiliary characters
are untwisted by the algebraic products (H10b).

The ordinary characters over \(\mathbb T\) are explicit.
Set \(U_{v,-1}=1\); prescribe them on units by (H5)
and on a uniformizer by
\[
        \psi_{vi}(g_v)\epsilon_\ell(g_v)^i
                                  =U_{vi}U_{v,i-1}^{-1}.
 \tag{H11}
\]
The ordinary operators are units in the complete local
algebra \(\mathbb T\).
Local reciprocity extends these data continuously:
the local multiplicative group is the product of its
units and the cyclic uniformizer group, whose map
extends to its profinite completion since
\(\mathbb T^\times\) is profinite.
At each separating point these are the actual ordered
Galois characters, by the two local equations and the
ordinary flag.

At the separating classical points the Galois local flags
used above do not follow from semisimplified compatibility
alone. Here is the additional check.
Distinct finite inertial characters make the smooth local
factor at \(\ell\) an irreducible principal series with
zero-monodromy parameter.
After a finite abelian extension killing those characters
its parameter is unramified.
Apply full compatibility and the crystalline assertion of
the cuspidal Galois theorem, after a compatible solvable
base change, and then restrict back.
This proves potential crystallinity and the asserted
finite inertial type, with zero monodromy.

Over that local extension, order the Hodge numbers
\(h_{\sigma0}<\cdots<h_{\sigma3}\) by the ordered ordinary
characters. The normalized unit equations for \(U_{vj}\)
say that the Newton number of the first \(j\) parameter
lines is
\(\sum_\sigma\sum_{i<j}h_{\sigma i}\).
For any \(j\)-dimensional filtered subspace this is the
minimum possible Hodge number.
Weak admissibility gives the opposite inequality, so
equality holds on these parameter subspaces.
They are weakly admissible and yield Galois subrepresentations.
Their successive rank-one quotients are the characters
of (H5), (H11).
Their distinct slopes make this flag canonical after the
chosen ordering, so it descends.
This supplies the actual ordinary flag, including its
unit and uniformizer values.
Both Galois-realization inputs used here apply at every
odd coefficient prime; in particular rank four at
\(\ell=3\) is included.
\subsubsection{The quotient and its arithmetic points}
Let \(B\) be the universal ordinary polarized determinant
algebra for the same residual link and local characters.

\begin{srproposition}\label{prop:sr-case-two-hecke-quotient}
There is a continuous surjection
\[
                             B\longrightarrow\mathbb T
 \tag{H12}
\]
compatible with \(\Lambda\), the Galois determinant and
the ordered ordinary characters.
Its kernel \(J_0\) is the actual automorphic ideal.
The quotient \(B/J_0\) is finite and torsion free over
\(\Lambda\), with dense classical arithmetic points.

The normalized Case II seed defines an arithmetic point
\(x\) of \(B\) with \(J_0\subseteq\ker x\).
Every regular arithmetic point of \(B/J_0\) is realized
by actual classical definite-unitary forms with the
specified semisimple Galois determinant.
If that Galois representation is irreducible, its
general-linear base change is cuspidal.
\end{srproposition}
\begin{proof}
At the separating classical points, the attached
representations have the specified unitary polarization
and ordered ordinary flags (H5), (H11).
At tame type places they have inertia (H2) and zero
monodromy, by (H3) and full local--global compatibility.
They are unramified outside the specified ramification set.
Thus the polarization, ordinary Cayley--Hamilton and
killed-inertia relations hold at every separating point.
They are polynomial determinant relations, so they hold
over \(\mathbb T\).
Equations (H7), (H6), Chebotarev and the integral identity
(H9) give the specified residual determinant and local
characters. Universality gives (H12).

Its image contains \(\Lambda\), all spherical operators
by (H7), and all ordinary operators by (H6).
These topologically generate \(\mathbb T\).
The image is closed by compactness of the source,
so (H12) is surjective.
Finiteness, torsion freeness and density follow from
(H8) and Lemma~\ref{lem:sr-unitary-classical-separation}.
If the residual module is zero, set \(J_0=B\).
The seed argument next shows that this alternative
does not occur for the actual seed link.

Use the seed's actual cuspidal automorphic representation
from Proposition~\ref{prop:sr-induced-case-two-seed},
with the normalization of
Proposition~\ref{prop:sr-seed-local-zero}.
Its prescribed regular weights, tame inertial characters
and nonexceptional local ordinary flag were checked there.
Its ordinary local Jacquet vector has exactly the unit
and uniformizer values (H5)--(H6).
The auxiliary \(v_0\) was chosen away from its ramification;
the seed and link-normalizing character are unramified
there. Its split local factor therefore has a vector fixed
by the first principal congruence subgroup.
At tame places (H3) supplies the required coefficient
vectors; at nonsplit places it is unramified.

The proved generic pure-inner-form descent described
after (H1) now gives a nonzero vector with this exact
weight, tame type and auxiliary level in the classical
definite-unitary space.
Classical control places its eigensystem in the
specialization of the constructed module.
Its residual values are (H7), (H6), so it belongs to \(M\).
It gives a character \(\mathbb T\to E_x\) whose composite
with (H12) equals \(x\):
the weight characters, Frobenius polynomials and ordered
ordinary characters agree, and they generate \(B\).
Hence \(J_0\subseteq\ker x\).

Finally \(M\) is a faithful finite \(\mathbb T\)-module.
Every prime of \(\mathbb T\) therefore belongs to its
support. Localize at a characteristic-zero point and
apply Nakayama to see that the corresponding fiber of
\(M\) is nonzero.
Classical control gives actual classical forms.
A common eigenvector for the commuting finite-dimensional
algebra in this fiber realizes the point.
Its attached semisimple Galois determinant equals the
given determinant by Frobenius compatibility and
Chebotarev. If this representation is irreducible,
(H10) has one summand and its length is one.
The base change is then cuspidal.
\end{proof}

The point property just proved uses actual automorphy,
the matched tame type, the nonexceptional ordered local
flag and unramifiedness at \(v_0\).
Those are the properties checked for the constructed seed.

\subsubsection{Speh strings and the required spectral bound}
The spectral assertion used at this endpoint is the
strict bound called \(\mathrm{PureW}\):
\[
               |\iota(\alpha)|<
                    \operatorname{Norm}(w)|\iota(\beta)|
 \tag{H13}
\]
for any two nonzero Frobenius roots, off finitely many
places. It is weaker than equality of their absolute
values and suffices for the later block exclusion.

\begin{srlemma}\label{lem:sr-case-two-speh-exclusion}
At an arithmetic point on the actual link quotient
\(B/J_0\), every length \(b_i\) in (H10) is one.
Bound (H13) follows.
Undoing the scalar link twist gives the
\(\mathrm{PureW}\) condition on the original rational
rank-four representation.
\end{srlemma}
\begin{proof}
The exact residual input is
\[
          \bar\chi_i\ne\bar\chi_j\bar\epsilon_\ell
                       \quad\hbox{on }G_L\quad(0\le i,j<4).
 \tag{H14}
\]
It includes \(i=j\).
On the tower this follows from
\[
 \forall w\in\mathbb Z^2\ \exists g\in G_L:\quad
       \overline{(\chi_0/\chi_1)^{w_1}
                 (\chi_0/\chi_2)^{w_2}}(g)
                              \ne\bar\epsilon_\ell(g).
 \tag{H15}
\]
The paired relation \(\chi_0\chi_3=\chi_1\chi_2\)
expresses every \(\chi_i/\chi_j\) by its root weight
\(w=w(i,j)\). A scalar link twist cancels.
The tower condition \(\mathrm{Ok.g2}\) is exactly
(H15). The subgroup of \(G_F\) fixing \(L\) is
canonically \(G_L\), so its witnesses belong to the
required group.

Choose stable lattices in the summands of (H10).
Their direct sum has residual semisimplification
\(\bigoplus_i\bar\chi_i\), independently of the lattice:
the characteristic polynomials agree with those of
the actual determinant, and characteristic-polynomial
Brauer--Nesbitt applies.
If \(b_i\ge2\), a nonzero Jordan--H\"older
constituent of the reduction of \(\rho_i'\) is
some \(\bar\chi_s\).
The next summand
\(\rho_i'\epsilon_\ell^{-1}\) makes
\(\bar\chi_s\bar\epsilon_\ell^{-1}\) another constituent,
say \(\bar\chi_t\).
This contradicts (H14). All \(b_i\) are therefore one.

The remaining constituents are unitary cuspidal
general-linear representations. Their unramified
local components are unitary and generic.
The generic part of the
\href{https://www.numdam.org/item/ASENS_1986_4_19_3_335_0/}
{Tadi\'c unitary classification} expresses them using
tempered factors and complementary pairs with
exponents \(s,-s\), \(0<s<1/2\).
At an unramified place all factors are spherical;
the Satake absolute values lie strictly between
\(q_w^{-1/2}\) and \(q_w^{1/2}\).
This holds for all constituents at once, so every
two parameters have ratio strictly less than \(q_w\).
The common rank-four Galois normalization multiplies
them all by the same scalar and preserves this ratio.
This proves (H13).

A scalar link twist also cancels from the ratio.
For a rational Frobenius power lying in \(G_L\),
let \(f_w\) be the residue degree of the place below it.
A rational exponent \(f\) is a multiple of \(f_w\);
its eigenvalues are the \(f/f_w\)-th powers of
the local Frobenius eigenvalues.
Thus the bound \(<q_w\) becomes \(<p^f\).
Exclude the finitely many ramified places and take the
open subgroup \(G_L\). This is the required
\(\mathrm{PureW}\) condition.
\end{proof}

\subsubsection{One module on the two links}
The shared arithmetic Taylor--Wiles calculation has an
exact role in this match.
On a selected nice characteristic-\(\ell\) curve, its
finite-level Selmer inputs are the first- and second-Selmer bounds
of Section~\ref{sec:sr-arithmetic-selmer-producers}, on the actual auxiliary
Galois groups with the actual ordinary and polarized
local conditions. Their constants precede precision
and auxiliary level.
That shared arithmetic calculation proves these estimates and
constructs the actual auxiliary groups and primes, under the nice-point
conditions displayed there. The link must supply those same conditions,
including its distinguished local \srcid{V0ok} condition.
The modules to which they apply are precisely the
ordinary function modules with lattice (H3),
character (H5), residual system (H7) and varying
split auxiliary local levels.
One module system is used along the prime chain.

The first finite-level hypothesis of
Proposition~\ref{prop:sr-one-patched-module}
is ordinary auxiliary projectivity and truncated
freeness.
Its second is given by the explicit old/new and
trace maps of
Proposition~\ref{prop:concrete-two-maps},
with the same normalized auxiliary discriminant.
For the third, use the actual Galois determinant
constructed above at each auxiliary level.
Its polarization, ordinary, inertia and auxiliary
relations are the relations of \(B\).
The explicit torus calculation gives nilpotence
exponent \(n!=24\) in rank four.
All maps act on the same free tame coefficient
lattice and commute with its tame action, so the
lattice introduces no new denominator.

Consequently the three finite-level hypotheses hold
for this actual system once the stated arithmetic
Selmer bounds and auxiliary construction are invoked.
Proposition~\ref{prop:tw-limit-coh} supplies the local
presentation estimates, and
Propositions~\ref{prop:sr-one-patched-module} and
\ref{prop:sr-completed-support-chain} give a common
module, its completed supports and their transport.
This is the rank-four application of those shared proofs.

Choose the two links with the same scalar character
\(\kappa\) and distinct auxiliary primes.
For their universal ordinary determinant algebras
\(B,B'\), take the actual kernels \(J_0,J'_0\)
in (H12). On their common paired space \(X\), put
\[
                        J=\phi(J_0)X+\phi'(J'_0)X.
 \tag{H16}
\]
The finite weight-zero property follows from
finiteness of \(B/J_0\) over \(\Lambda\).
Indeed take a characteristic-\(\ell\) domain
quotient after setting the fixed inertial weight
characters to \(1\).
The image of \(\Lambda\) is finite, so the link
Hecke quotient is algebraic over a finite field.
Thus all characteristic coefficients on \(G_L\)
are algebraic constants.
Every element of \(G_F\) has square in \(G_L\).
Its eigenvalues and characteristic coefficients
are algebraic constants as well.
The complete local domain has the original finite
residue field. Reduction is injective on every finite
subfield, so any algebraic constant in it is the unique
corresponding element of that residue field.
All characteristic coefficients thus equal their
residual values.
These coefficient differences and the inertial weight
differences topologically generate the maximal ideal
of the paired trace algebra. Formal Nakayama, using
Noetherianity, makes that maximal ideal zero in the
domain quotient. Every such prime is therefore maximal.
A power of the maximal ideal is consequently contained
in the fixed-weight ideal, as required in
Proposition~\ref{prop:sr-four-core}.

On each nice curve the common module, the local
presentations and
Proposition~\ref{prop:sr-four-curve-genericity}
give the curve-local support assertion.
At the seed,
Proposition~\ref{prop:sr-case-two-hecke-quotient}
gives \(J_0x=J'_0x=0\), and
Proposition~\ref{prop:sr-seed-local-zero}
gives local vanishing of their two mapped ideals.
The residual chain and
Proposition~\ref{prop:sr-four-core} then transport
this support to the original paired point.
The characteristic-\(\ell\) exclusion is
Lemma~\ref{lem:sr-four-characteristic-ell};
the characteristic-zero spectral assertion is
Lemma~\ref{lem:sr-case-two-speh-exclusion}.

At a regular irreducible arithmetic point, classical
control gives an automorphic representation with
the prescribed local factors away from the
distinguished auxiliary place.
Both links describe the same global general-linear
representation: their scalar twist is the same,
and their Frobenius polynomials agree away from
their two auxiliary places.
Strong multiplicity one identifies their cuspidal
representations. At any one finite place at
least one link has no auxiliary exception.
Its local--global comparison therefore gives
the desired local statement there, including
at the other link's auxiliary place.
The level introduced for smallness has thereby
been removed from the actual endpoint.

This identifies the Case II Hecke and support
objects with the shared ordinary construction.
The arithmetic production of those Selmer bounds and
of the actual auxiliary groups and prime-selection limit is supplied
by Section~\ref{sec:sr-arithmetic-selmer-producers}; the matching nice-point
and common-space conditions remain the ones explicitly used above.
The further passage from automorphy over \(L\)
to the source's \(\mathrm{IsAutQA}\), including
its \(\mathrm{WitA}\) local level data over an
imaginary quadratic field, is the separate
solvable-descent and local-compatibility step.
It is not inferred solely from the existence
of a Hecke eigensystem.

% Exact frozen-source correspondence:
% DenseThresholdCases/HdEx, TTypeEx, TThetaEx, TResEx,
% UnitaryHeckeData, TDenseU, TCoreSep, TCoreDet, OrdCoreH,
% ExistsHeckeIdeal, TJRing, TJDense, THidaU, TLevelU, THasAutOf;
% PointTheoremBaseObject, BirthPm, TPmFin, PmTwo,
% Gsp4PromodularityCaseTwo, VWinPm, PmAutO, DnPmDown.
% NoOmLift and MemMax in
% WeakShimuraTheory/CharacteristicPolynomialRoots.
